\section{Discussion}

\subsection{Interpreting the Null (and Negative) Result}

The curated-corpus hypothesis predicts $R_{\text{OLMo}} > R_{\text{Pythia}}$. We observe the opposite. Three explanations, with assessed plausibility, follow.

\subsubsection{Explanation 1: Architecture Difference (HIGH plausibility)}

OLMo-7B and Pythia-6.9B differ in activation function (SwiGLU vs.\ GELU), tokenizer vocabulary size, positional embedding scheme, and optimizer configuration. These differences are well-documented~\cite{Groeneveld2024OLMo,Biderman2023Pythia} and are known to affect downstream benchmark performance independently of training data. MMLU performance is sensitive to tokenizer vocabulary and few-shot prompt formatting, both of which differ between the two models. The observed MMLU gap (Pythia: 0.2588 vs.\ OLMo: 0.2464) may reflect Pythia's architectural advantages on MMLU-formatted prompts rather than corpus quality differences. This explanation is highly plausible and cannot be ruled out without architecture-matched experiments.

\subsubsection{Explanation 2: Training Scale Insufficient at 300B Tokens (HIGH plausibility)}

Corpus quality effects on knowledge acquisition may require more training tokens to manifest. At 300B tokens, both models may be in a regime where architecture-driven learning efficiency dominates corpus quality effects. Scaling studies~\cite{Muennighoff2023Scaling} suggest that data quality effects interact with compute budget in complex ways; the 300B token scale may be below the threshold at which Dolma's academic content inclusions provide measurable advantage. Longitudinal trajectory analysis across multiple checkpoints (e.g., 100B, 200B, 300B, 400B tokens) would test whether the gap narrows or reverses at higher scales.

\subsubsection{Explanation 3: The Pile MMLU Contamination (LOW-MEDIUM plausibility)}

The Pile contains web-crawled content that may overlap with MMLU questions or their sources. If Pythia-6.9B has seen MMLU-adjacent content during training, its higher MMLU score may reflect contamination rather than genuine knowledge acquisition. This explanation is plausible but lower priority than the architecture explanation, because: (a) contamination effects on 5-shot MMLU are typically modest; (b) Dolma also contains web-crawled content; and (c) the per-subject heatmap does not show the uniform inflation pattern characteristic of systematic contamination. Definitive assessment requires n-gram overlap analysis~\cite{shi2024detecting} between The Pile and MMLU questions.

\subsubsection{The Structural Insight: HellaSwag Denominator Saturation}

Regardless of which explanation accounts for the MMLU gap, the HellaSwag convergence finding is structurally important and not explained by the above alternatives. Both models achieving exactly 0.4580 on HellaSwag --- despite substantial architectural and corpus differences --- is consistent with a saturation phenomenon: at $\sim$300B tokens, 6--8B parameter models have largely exhausted the commonsense signal in web-sourced data, and HellaSwag becomes insensitive to further corpus quality variation. This is supported by prior work showing near-zero sensitivity of HellaSwag to deduplication interventions~\cite{Penedo2023RefinedWeb}. The implication for metric design is direct: ratio metrics using HellaSwag as a denominator lose discriminative power at this scale and should be validated for denominator sensitivity before adoption in corpus quality comparisons.

\subsection{Limitations}

\textbf{L1: Architecture confound.} OLMo-7B and Pythia-6.9B are not architecture-matched. Performance differences cannot be attributed to corpus quality alone.
\textit{Why acceptable:} The study's contribution is to demonstrate that the confound matters and to quantify the magnitude of the observed difference, motivating architecture-matched follow-up.
\textit{Future work:} Repeat comparison using Pythia-12B (closer parameter count) or OLMo models trained on The Pile as an ablation.

\textbf{L2: 500-sample fast evaluation.} The \texttt{--limit 500} flag subsamples benchmarks, introducing sampling variance.
\textit{Why acceptable:} Bootstrap CI accounts for this variance; the HellaSwag convergence finding should be verified with full evaluation.
\textit{Future work:} Full benchmark evaluation (estimated 4--6 hours per model on H100) to confirm HellaSwag exact convergence.

\textbf{L3: Single training scale.} Results are specific to $\sim$300B tokens and may not generalize to other points on the training curve.
\textit{Why acceptable:} 300B tokens is a well-documented, practically relevant scale for both models, and both provide public checkpoints.
\textit{Future work:} Multi-checkpoint trajectory analysis to identify where (if anywhere) corpus quality effects emerge.

\textbf{L4: No contamination analysis.} We do not perform n-gram overlap analysis between corpora and benchmarks.
\textit{Why acceptable:} Contamination is a lower-plausibility explanation than architecture effects, and the per-subject heatmap does not show uniform inflation.
\textit{Future work:} Apply \texttt{min-k\%} or n-gram overlap methods~\cite{shi2024detecting} to both corpora.

\subsection{Implications for Metric Design}

The results motivate three concrete recommendations for future corpus quality studies:

\begin{enumerate}
    \item \textbf{Verify denominator sensitivity before adopting ratio metrics.} A ratio metric in which the denominator does not vary across model variants is a MMLU proxy, not a generalization balance metric. Confirm that HellaSwag (or any denominator benchmark) shows meaningful variance across the models being compared before using it as a denominator.
    \item \textbf{Use architecture-matched baselines or trajectory analysis.} Single-point comparisons between architecturally distinct models cannot isolate corpus quality effects. Architecture-matched designs (same model trained on different corpora) or temporal trajectory analysis (ratio evolution across training checkpoints) are the minimum required designs for corpus quality attribution.
    \item \textbf{Report effect sizes and CIs, not just point estimates.} The large Cohen's $d$ ($-2.732$) in this study is inflated by small per-subject variance at 500-sample evaluation. Full evaluation and effect size reporting with appropriate uncertainty quantification are necessary for valid corpus quality comparisons.
\end{enumerate}

\subsection{Broader Impact}

This work contributes to the reproducibility infrastructure for large language model research by providing a documented, publicly reproducible evaluation pipeline using open checkpoints. The null result is itself a contribution: the field invests heavily in corpus curation based on the assumption that curation quality translates to generalization gains, and our results suggest this assumption requires more careful empirical validation at controlled scales. Publishing null results with full statistical characterization supports more accurate calibration of research investment in data curation.
